% BEGIN REV09_Q_COMPLEJO
\subsection{Iteración del lector de razones y recuperación de fase}
\label{corpus9:iv:q}

El lector \(R(x,y)=x/y^2\), aplicado a las representaciones positivas
de clausura y autoescala, admite la composición orientada
\[
 \mathcal Q_{\mathrm{rat}}(x,y)
   =\bigl(R(x,y),R(y,x)\bigr)
   =(x/y^2,y/x^2),\qquad x,y>0.
\]
Esta composición recibe las coordenadas ya producidas por
APP--TRIT--TPK. Su estudio conserva tanto el producto \(P=xy\)
como la razón orientada \(O=x/y\), cuyas transformaciones son
\[
 P\circ\mathcal Q_{\mathrm{rat}}=P^{-1},\qquad
 O\circ\mathcal Q_{\mathrm{rat}}=O^3.
\]
La notación \(\mathcal Q_{\mathrm{rat}}\) distingue esta aplicación
de la extensión a intervalos del lector escalar.

\begin{proposition}[Iteración positiva y prolongación compleja]
\label{corpus9:iv:q-iteracion}
Para \(n\geq0\), sean
\[
 a_n=\frac{(-1)^n+3^n}{2},\qquad
 b_n=\frac{(-1)^n-3^n}{2}.
\]
Entonces
\begin{equation}
 \mathcal Q_{\mathrm{rat}}^n(x,y)
   =(x^{a_n}y^{b_n},x^{b_n}y^{a_n}).
 \label{corpus9:iv:q-potencias}
\end{equation}
La misma expresión de Laurent define una aplicación sobre
\((\mathbb C^\times)^2\). Para \(n\geq1\), esta aplicación es un
recubrimiento no ramificado de grado \(3^n\), con núcleo
\begin{equation}
 \ker\mathcal Q_{\mathrm{rat}}^n
   =\{(\zeta,\zeta^{-1}):\zeta^{3^n}=1\}.
 \label{corpus9:iv:q-nucleo}
\end{equation}
Su restricción al cuadrante positivo tiene una preimagen positiva
única para cada pareja positiva.
\end{proposition}
\begin{proof}
En coordenadas logarítmicas positivas, la matriz es
\[
 B=\begin{pmatrix}1&-2\\-2&1\end{pmatrix}.
\]
Los vectores \((1,1)\) y \((1,-1)\) tienen valores propios
\(-1\) y \(3\). La descomposición en esas dos direcciones da
\(B^n=\left(\begin{smallmatrix}a_n&b_n\\b_n&a_n\end{smallmatrix}\right)\)
y prueba \eqref{corpus9:iv:q-potencias}. Como \(a_n,b_n\) son
enteros, la identidad se prolonga como identidad de Laurent.

Escribamos \(m=3^n\), \(\epsilon=(-1)^n\). Para una imagen
\((u,v)\in(\mathbb C^\times)^2\), el producto de una preimagen
queda fijado por \(p=xy=(uv)^\epsilon\). Las ecuaciones equivalen a
\[
 x^m=u\,p^{-b_n},\qquad y=p/x.
\]
La primera tiene exactamente \(m\) raíces distintas. La segunda
determina \(y\); además, el producto de las dos imágenes es
\(p^\epsilon=uv\), de modo que la segunda ecuación original se
cumple. Toda fibra contiene exactamente \(m\) puntos. Para
\((u,v)=(1,1)\), resulta \eqref{corpus9:iv:q-nucleo}; las otras
fibras son sus trasladados en el grupo multiplicativo.

El diferencial satisface
\[
 \det D\mathcal Q_{\mathrm{rat}}^n
   =(-3)^n(xy)^{\epsilon-1}\ne0.
\]
La descripción local por raíces de un número no nulo proporciona
las \(m\) ramas locales y prueba la afirmación de recubrimiento.
En el cuadrante positivo sólo una de esas raíces es positiva.
En particular, para un paso,
\[
 x=(uv^2)^{-1/3},\qquad y=(u^2v)^{-1/3},
\]
con las raíces reales positivas.
\end{proof}

En la prolongación compleja, la condición \(xy=(uv)^{-1}\)
acopla las dos raíces cúbicas del inversor de un paso. Las
coordenadas \((P,O)\) identifican además
\((x,y)\) y \((-x,-y)\); su aplicación tiene núcleo
\(\{(1,1),(-1,-1)\}\). La prueba anterior conserva \((x,y)\)
y, con ello, la distinción entre esa identificación de coordenadas
y las fibras de \(\mathcal Q_{\mathrm{rat}}^n\).

\begin{proposition}[Reconstrucción por el registro ternario]
\label{corpus9:iv:q-acarreo}
Sobre el subgrupo
\[
 (x,y)=(e^{2\pi i\theta},e^{-2\pi i\theta}),
\]
la aplicación induce \(\theta\mapsto3\theta\pmod1\).
Fijados los representantes \(\theta_j\in[0,1)\), sean
\[
 d_j=\lfloor3\theta_j\rfloor,\qquad
 \theta_{j+1}=3\theta_j-d_j,\qquad
 C_n=\sum_{j=0}^{n-1}d_j3^{n-1-j}.
\]
La recuperación exacta es
\begin{equation}
 \theta_0=\frac{\theta_n+C_n}{3^n}.
 \label{corpus9:iv:q-recuperacion}
\end{equation}
\end{proposition}
\begin{proof}
La sustitución en la expresión de Laurent da la triplicación de
fase. La recurrencia \(C_{n+1}=3C_n+d_n\), con \(C_0=0\),
da por inducción \(\theta_n=3^n\theta_0-C_n\). Esto prueba
\eqref{corpus9:iv:q-recuperacion}. Cada \(d_j\) pertenece a
\(\{0,1,2\}\), por lo que \(n\) dígitos ternarios conservan
la elección entre las \(3^n\) preimágenes de la fase reducida.
Las ramas \(B_d(\theta)=(\theta+d)/3\) satisfacen
\[
 B_a\circ B_b(\theta)=\frac{\theta+b+3a}{9}.
\]
Por tanto, dos pasos conservan el dígito ordenado \(3d_0+d_1\).
\end{proof}

El producto retorna después de dos iteraciones, mientras la razón
orientada se eleva a la novena potencia. La recuperación demostrada
tiene por dominio este lector y su registro de fase. Su inserción
en una representación del estado enriquecido completo conserva
adicionalmente los mapas de hojas, orientación, memoria y
supervivencia propios de esa representación.
% END REV09_Q_COMPLEJO
